\section{Assessment}
\label{sec:assessment}

We assess how well Kairo localizes annotated changes in complete event logs. The event logs, detector configurations, and scoring protocol provide a common basis for comparison with implementations in \softwarepm{}. \Cref{sec:assessment-results} addresses \ref{rq:comparison} with results pooled across logs. \Cref{sec:assessment-scenarios} answers \ref{rq:scenarios} through selected change scenarios. \Cref{sec:assessment-sensitivity} addresses \ref{rq:margin} by repeating the scoring with different matching margins.

\subsection{Event Logs}
\label{sec:assessment-logs}

We selected 23 event logs with 156,383 cases, 1,086,668 events, and 37 annotated drift episodes. \autoref{tab:assessment-logs} lists all inputs. The selection covers several kinds of change and uses each selected log in full. The Padella collection injects changes into logs from real processes~\cite{padella2026dataset}. Our sample uses eight files from five base processes, including two controls with no injected drift. The nine synthetic Rheon logs from the thesis include sudden and gradual changes to activity order, assignments, workload, waiting time, and case composition~\cite{kairo-thesis}. The six historical logs from the archive of Adams et~al. contain synthetic control-flow changes~\cite{adams2023experimental}. They include a Ceravolo scenario with 15\% trace noise~\cite{tavares2019synthetic}.

\begin{table}[H]
\caption{Event logs used in the comparison. Simultaneous annotations with identical time bounds count as one episode.}
\label{tab:assessment-logs}
\centering
\small
\begin{adjustwidth}{-\extralength}{0cm}
\begin{tabularx}{\fulllength}{@{}lrrr>{\raggedright\arraybackslash}X@{}}
\toprule
\textbf{Log} & \textbf{Cases} & \textbf{Events} & \textbf{Episodes} & \textbf{Annotated changes} \\
\midrule
Bose & 6,000 & 58,838 & 4 & Sudden control-flow changes \\
Ceravolo RIO & 1,000 & 11,395 & 1 & Sudden control-flow change \\
Ceravolo RIO (15\% noise) & 1,000 & 10,554 & 1 & Sudden control-flow change \\
Ceravolo RP & 1,000 & 11,009 & 1 & Sudden control-flow change \\
Ostovar Loop & 2,999 & 100,471 & 2 & Sudden changes involving a loop \\
Ostovar Parallel to Sequence & 2,999 & 112,011 & 2 & Sudden changes involving parallel and sequential execution \\
\midrule
Padella 04 & 8,049 & 51,059 & 1 & Gradual speedup in loan applications \\
Padella 05 & 10,500 & 52,308 & 1 & Sudden control-flow change in domestic declarations \\
Padella 06 & 32,429 & 212,721 & 1 & Sudden speedup in bank account closures \\
Padella 07 & 6,886 & 36,796 & 0 & Request for payment control \\
Padella 08 & 6,886 & 36,796 & 1 & Sudden control-flow change in requests for payment \\
Padella 10 & 10,500 & 56,437 & 1 & Sudden slowdown in domestic declarations \\
Padella 12 & 8,049 & 51,059 & 0 & Loan application control \\
Padella 15 & 8,000 & 39,493 & 1 & Gradual control-flow change in purchase orders \\
\midrule
Rheon Single 03 & 3,254 & 12,992 & 1 & Sudden control-flow change \\
Rheon Single 11 & 8,896 & 24,543 & 1 & Sudden resource reassignment \\
Rheon Single 15 & 9,042 & 72,684 & 1 & Gradual change in waiting time \\
Rheon Multi 02 & 3,161 & 20,514 & 4 & Sudden control-flow changes \\
Rheon Multi 03 & 9,523 & 23,852 & 3 & Sudden changes in the regional case mix \\
Rheon Multi 07 & 3,132 & 15,660 & 3 & Changes in workload, resource pool size, and assignments \\
Rheon Mixed 06 & 4,796 & 23,980 & 1 & Coincident workload and waiting-time changes \\
Rheon Mixed 07 & 2,860 & 14,300 & 2 & Changes in workload and case arrival rates \\
Rheon Mixed 13 & 5,422 & 37,196 & 4 & Changes in control flow, workload, and case arrival rates \\
\midrule
\textbf{Total} & \textbf{156,383} & \textbf{1,086,668} & \textbf{37} & \\
\bottomrule
\end{tabularx}
\end{adjustwidth}
\end{table}

We expressed all annotations in calendar time. For Rheon we used generator timestamps. For Padella we used the start times of the cases identified in the published annotations. For Bose and Ostovar we used the start of the first post-change trace in the original archive order. The Ceravolo archive stores numeric case identifiers in lexicographic order, so we used the start of case 500, the first case in the changed half by numeric generator identifier. We preserved physical trace boundaries when the Bose archive reused a case identifier.

\subsection{Detectors and Configurations}
\label{sec:assessment-config}

The comparison covers nine Kairo combinations of perspective and clustering method and 18 \softwarepm{} implementations (\autoref{tab:assessment-methods}). We call a specific configuration of an implementation a \emph{profile}. Core profile settings used observable input properties, such as log length, event and case volume, and available fields. The annotated change times were used for scoring after detection. The core runs used Kairo 0.1.0 at revision \texttt{6b1bd462} and \softwarepm{} at revision \texttt{79b6ce1b}, with revision \texttt{765435bc} for the corrected Bose detector. The benchmark archive records the complete settings for every run. Coverage depended on data and computational requirements. The \softwarepm{} resource-assignment profile ran only when the log recorded resources. Kairo recorded 30 unsupported runs due to missing resources and eight because the number of distinct intra-case feature rows exceeded the DBSCAN limit of 10,000. In total, 616 \softwarepm{} runs and 376 Kairo runs completed.

\begin{table}[H]
\caption{Implementations and profiles included in the core comparison. DFG denotes a directly-follows graph. PELT denotes pruned exact linear time, ADWIN denotes adaptive windowing, and MMD denotes maximum mean discrepancy.}
\label{tab:assessment-methods}
\centering
\small
\begin{adjustwidth}{-\extralength}{0cm}
\begin{tabularx}{\fulllength}{@{}l>{\raggedright\arraybackslash}p{0.19\fulllength}>{\raggedright\arraybackslash}X@{}}
\toprule
\textbf{Library} & \textbf{Implementation} & \textbf{Signal and tested profiles} \\
\midrule
Kairo & State distributions & Intra-case, resource, and inter-case states with K-means, a self-organizing map (SOM), or DBSCAN. Each combination uses previous-window and recent-window profiles. \\
\midrule
\softwarepm{} & Bose relation count & Permutation test on counts of outgoing eventually-follows relations. \\
& DFG frequency & Adjacent-window comparison of relation frequencies using permutations or a wider likelihood-ratio test. \\
& Event relations & Fixed or adaptive windows of event-level control-flow relations. \\
& Local completeness & New or disappearing directly-follows relations relative to a locally complete window. \\
& Relation invariance & Boundaries of stable periods for activity relations. \\
& Graph metrics & PELT on measures of successive DFGs. \\
& Run distribution & Fixed or adaptive windows of partially ordered execution runs. \\
& Trace variant frequency & Permutation tests on exact activity sequences with regular and wide windows. \\
& Earth Mover distance & Window comparisons with activity sequences alone or with timing features. \\
& Calendar indicators & PELT on basic or enriched indicator sets, and ADWIN on basic indicators. \\
& CV-based indicators & PELT on enriched indicators with smoothing selected using the coefficient of variation. \\
& Conformance & Alignment fitness monitored with fixed-model Page-Hinkley or ADWIN with relearning. \\
& Conformance profile & PELT on windowed alignment fitness and precision. \\
& Declare profiles & PELT on temporal constraint profiles. \\
& Version clustering & Clustering and segmentation of window similarity profiles based on DFGs. \\
& Cluster frequency & Comparison of trace cluster frequencies after training K-means on an initial set of traces. \\
& Multivariate features & Permutation MMD on flow and timing features or on activity-resource assignments. \\
& Supervised similarity & A patch classifier on window similarity images trained on independent generated logs. \\
\bottomrule
\end{tabularx}
\end{adjustwidth}
\end{table}

\paragraph{Kairo profiles.}
The state models were fitted to each complete log. The benchmark applied min-max scaling to intra-case vectors and standard score scaling to resource and inter-case vectors, followed by principal component analysis with up to five components when applicable. It used observed fields and feature size limits to choose the vectors. The base signal window targeted one fiftieth of the log's completion-time span, rounded to an available duration. Resource and inter-case feature windows were chosen from lengths that divide the base signal window, taking event or case density into account. The previous-window profile used Jensen-Shannon divergence and the 90th percentile of eligible scores as its threshold. The recent-window profile used total variation relative to the mean of the preceding three populated windows, twice the base signal window width, and the 95th percentile. Both profiles required full coverage and at least 25\% of the median number of state observations. \Cref{def:assessment-alarms} specifies how these settings produce retrospective alarms from the scores.

\begin{Definition}[Retrospective alarm rule]\label{def:assessment-alarms}
Let $t_{\min}$ and $t_{\max}$ be the first and last completion timestamps in the input log, and let $\widetilde N$ be the median of the populated-window counts $N_j$ from \cref{def:state-distributions}. For a chosen support fraction $\rho$, the set of eligible indices is
\[
  J=\{j:\lambda_j\text{ is finite},\ V_j\subseteq[t_{\min},t_{\max}),\
               N_j\geq\rho\widetilde N\}.
\]
If $J$ is nonempty, compute the empirical quantile at the chosen level from the sequence $(\lambda_j)_{j\in J}$, including repeated scores. The threshold $\theta$ is that quantile with a floor of $10^{-12}$. A window qualifies if $j\in J$ and $\lambda_j\geq\theta$. A maximal run of qualifying windows whose calendar boundaries touch produces an alarm at the start of its first window. An empty $J$ produces no alarms.
\end{Definition}

\paragraph{\softwarepm{} profiles.}
We used the recorded \softwarepm{} source revisions and retained each implementation's own alarms. Trace window sizes generally scaled with case count, and calendar indicators used bins sized by log duration. When usable service timestamps were missing, timing profiles used gaps between completions. Version clustering used 64 trace windows with a minimum cluster size of four and a segmentation penalty of three. The supervised similarity profile used a local classifier fitted to 24 independent synthetic control-flow logs. For the conformance profile, precision was calculated with the memory-efficient exact Dijkstra solver. On Ostovar Parallel to Sequence, exact alignment took too long with the model first mined from the initial traces. The conformance profiles used a sequential reference based on the medoid of the first 30 completed traces. Detection began after that training prefix, and the adaptive profile could relearn from one completed case following an alarm.

Before scoring the Bose detector, we corrected two implementation issues. The detector had used case identifiers to order traces and counted incoming eventually-follows relations. We ordered cases by their first event and used outgoing counts, as specified by the Relation Type Count definition~\cite{bose2011handling}. A check with permuted case identifiers confirmed that chronology was preserved. On the historical Bose log, the relation-count correction preserved three of four matches at the primary tolerance and reduced unmatched alarms from 17 to 12.

\subsection{Scoring Protocol}
\label{sec:assessment-scoring}

We used a common calendar-time margin equal to 2.5\% of the log's observation span, from the earliest recorded start or completion to the last completion. We repeated scoring with margins of 1\% and 5\%. Coincident annotations with identical bounds form one episode. \Cref{def:assessment-matching} applies the same localization rule to every detector, regardless of the type of change.

\begin{Definition}[Temporal matching]\label{def:assessment-matching}
An annotated episode $i$ has bounds $[l_i,r_i]$, with $l_i=r_i$ for a point change. The distance from an alarm timestamp $t$ to this episode is
\[
  d(t,[l_i,r_i])=\max\{l_i-t,\ 0,\ t-r_i\}.
\]
For a given margin $\gamma$, an alarm and an episode can be paired when their distance is at most $\gamma$. Among one-to-one matchings of feasible pairs, we maximize the number of matches and then minimize the sum of their distances.
\end{Definition}

Each match contributes a true positive ($TP$), each unmatched alarm a false positive ($FP$), and each unmatched episode a false negative ($FN$). Unmatched alarms include repeat alarms and changes that may have occurred outside the annotated injections. We sum $TP$, $FP$, and $FN$ across logs before computing micro precision $TP/(TP+FP)$, micro recall $TP/(TP+FN)$, and micro F1 $2TP/(2TP+FP+FN)$. The no-injection controls therefore contribute unmatched alarms to the pooled counts. We score the estimated change times. Measuring reporting delay would require the confirmation times, which the benchmark records where available.

\subsection{Comparative Performance Across Logs}
\label{sec:assessment-results}

For \ref{rq:comparison}, \autoref{tab:assessment-main-results} shows the three highest micro F1 profiles from each library among the profiles that completed on all 23 logs. At the primary margin, Kairo's previous-window intra-case K-means profile matched 21 of 37 episodes, while \softwarepm{} version clustering matched 15. Kairo raised 72 unmatched alarms and version clustering raised 10. Nine of Kairo's unmatched alarms and two of version clustering's came from the control logs. Version clustering therefore led in precision and micro F1, and the Kairo profile led in recall. Both profiles matched five episodes in the historical logs. On Padella, Kairo matched four with 27 unmatched alarms and version clustering matched three with seven unmatched alarms. On Rheon, Kairo matched 12 with 30 unmatched alarms and version clustering matched seven without an unmatched alarm.

\begin{table}[htbp]
\caption{Highest micro F1 profiles with full log coverage at the primary tolerance of 2.5\%. Rows are ranked within each library.}
\label{tab:assessment-main-results}
\centering
\small
\begin{tabularx}{\textwidth}{@{}l>{\raggedright\arraybackslash}Xrrrrrr@{}}
\toprule
\textbf{Library} & \textbf{Profile} & \textbf{TP} & \textbf{FP} & \textbf{FN} & \textbf{Precision} & \textbf{Recall} & \textbf{F1} \\
\midrule
Kairo & Intra-case K-means, previous & 21 & 72 & 16 & 0.226 & 0.568 & \textbf{0.323} \\
& Inter-case SOM, previous & 22 & 80 & 15 & 0.216 & 0.595 & 0.317 \\
& Intra-case SOM, previous & 20 & 73 & 17 & 0.215 & 0.541 & 0.308 \\
\softwarepm{} & Version clustering & 15 & 10 & 22 & 0.600 & 0.405 & \textbf{0.484} \\
& Declare profiles, PELT & 10 & 9 & 27 & 0.526 & 0.270 & 0.357 \\
& Trace variant frequency, wide & 20 & 66 & 17 & 0.233 & 0.541 & 0.325 \\
\bottomrule
\end{tabularx}
\end{table}

The coverage restriction matters for interpreting the pooled comparison. Kairo's recent-window resource K-means profile reached micro F1 of 0.370 on its 18 supported logs. On the same 18 logs, version clustering reached 0.440. We report such profiles separately because the omitted logs alter the set of episodes and opportunities for unmatched alarms.

\subsection{Detection by Change Scenario}
\label{sec:assessment-scenarios}

For \ref{rq:scenarios}, we compare profiles in selected scenarios at the primary matching margin. The examples show which annotated changes the perspective-specific state signals can capture under the tested settings.

\paragraph{Control flow.}
On the three Padella control-flow logs, Kairo's recent-window intra-case SOM and \softwarepm{} version clustering each matched all three injections and each raised two unmatched alarms. On Rheon Multi 02, Kairo's previous-window intra-case SOM matched all four sudden control-flow changes and raised one unmatched alarm. The wide trace variant frequency profile matched all four with no unmatched alarms. Both DFG frequency profiles matched all three Ceravolo injections with no unmatched alarms, including the log with trace noise.

\paragraph{Assignments and timing.}
All three recent-window Kairo resource profiles matched the sudden reassignment in Rheon Single 11 with no unmatched alarms. The \softwarepm{} resource MMD and enriched indicator PELT profiles also matched it and each raised one unmatched alarm. Kairo's recent-window inter-case SOM matched the gradual waiting-time change in Rheon Single 15 with no unmatched alarms. The trained trace cluster frequency profile matched it and raised one unmatched alarm. \autoref{fig:assessment-rheon-signals} pairs the Kairo resource and inter-case SOM signals for these two changes. On the three Padella timing logs, previous-window resource SOM and fixed-model Page-Hinkley conformance each matched two injections and raised nine unmatched alarms. Both missed the speedup in bank account closures.

\begin{figure}[H]
  \centering
  % Data for the side-by-side assessment panels.
% The source paths below are relative to assessments/concept-drift/results/.
% Calendar coordinates are normalized from the first window start to the last window end.
% Scores keep their original scale, with coordinates rounded to six decimal places.

% rheon_single_11/dataset.json and rheon_single_11/kairo/resource-som/recent.json.
\def\RheonReassignmentPath{%
  \draw[rheon/series] plot[mark=*,mark size=1.05pt] coordinates {
    (0.093750,0.240741) (0.125000,0.092593) (0.156250,0.074074) (0.187500,0.148148)
    (0.218750,0.074074) (0.250000,0.092593) (0.281250,0.240741) (0.312500,0.111111)
    (0.343750,0.259259) (0.375000,0.259259) (0.406250,0.777778) (0.437500,0.759259)
    (0.468750,0.444444) (0.500000,0.314815) (0.531250,0.222222) (0.562500,0.111111)
    (0.593750,0.111111) (0.625000,0.277778) (0.656250,0.111111) (0.687500,0.240741)
    (0.718750,0.074074) (0.750000,0.129630) (0.781250,0.240741) (0.812500,0.055556)
    (0.843750,0.222222) (0.875000,0.148148) (0.906250,0.203704) (0.937500,0.092593)
    (0.968750,0.444444)
  };%
}
\def\RheonReassignmentEdge{(0.968750,0.444444)}
\def\RheonReassignmentAlarm{(0.406250,0.777778)}
\def\RheonReassignmentThreshold{0.649074}
\def\RheonReassignmentOnset{0.414687}
\def\RheonReassignmentTicks{0.026042/Apr 2020, 0.343750/Jun, 0.661458/Aug, 0.979167/Oct}

% rheon_single_15/dataset.json and rheon_single_15/kairo/inter-som/recent.json.
\def\RheonWaitingPath{%
  \draw[rheon/series] plot[mark=*,mark size=1.05pt] coordinates {
    (0.100000,0.392857) (0.133333,0.273810) (0.166667,0.464286) (0.200000,0.190476)
    (0.233333,0.047619) (0.266667,0.107143) (0.300000,0.095238) (0.333333,0.071429)
    (0.366667,0.166667) (0.400000,0.607143) (0.433333,0.869048) (0.466667,0.785714)
    (0.500000,0.511905) (0.533333,0.369048) (0.566667,0.285714) (0.600000,0.142857)
    (0.633333,0.047619) (0.666667,0.154762) (0.700000,0.642857) (0.733333,0.369048)
    (0.766667,0.250000) (0.800000,0.535714) (0.833333,0.190476) (0.866667,0.345238)
    (0.900000,0.750000) (0.933333,0.761905) (0.966667,0.531746)
  };%
}
\def\RheonWaitingEdge{(0.966667,0.531746)}
\def\RheonWaitingAlarm{(0.433333,0.869048)}
\def\RheonWaitingThreshold{0.779762}
\def\RheonWaitingOnset{0.364286}
\def\RheonWaitingEnd{0.468571}
\def\RheonWaitingTicks{0.016667/2015, 0.233929/2016, 0.451786/2017, 0.669048/2018, 0.886310/2019}

  \begin{tikzpicture}[kairo/base,x=1cm,y=1cm,
      rheon/grid/.style={draw=KairoLine,line width=.35pt},
      rheon/axis/.style={draw=KairoMuted,line width=.5pt},
      rheon/series/.style={draw=KairoBlue,line width=.8pt,line join=round},
      rheon/threshold/.style={draw=KairoMuted,line width=.6pt,dash pattern=on 3pt off 2pt},
      rheon/truth/.style={draw=KairoAmber,line width=.65pt},
      rheon/alarm/.style={rectangle,rotate=45,inner sep=0pt,minimum size=5pt,
                          fill=KairoTeal,draw=white,line width=.45pt},
      rheon/tick/.style={inner sep=0pt,text=KairoMuted}]
    \node[anchor=west,font=\small\bfseries,inner sep=0pt] at (.83,2.61)
      {(a) Resource reassignment};
    \node[anchor=west,font=\small\bfseries,inner sep=0pt] at (7.39,2.61)
      {(b) Gradual waiting-time change};
    \node[rotate=90,inner sep=0pt,text=KairoMuted] at (.02,1.14) {Total variation};

    \begin{scope}[shift={(.83,0)},x=5.89cm,y=2.28cm]
      \fill[KairoPaper] (0,0) rectangle (1,1);
      \draw[rheon/grid] (0,.5) -- (1,.5) (0,1) -- (1,1);
      \draw[rheon/truth] (\RheonReassignmentOnset,0) -- (\RheonReassignmentOnset,1);
      \draw[rheon/threshold] (0,\RheonReassignmentThreshold) -- (1,\RheonReassignmentThreshold);
      \RheonReassignmentPath
      \filldraw[draw=KairoBlue,fill=white,line width=.55pt] \RheonReassignmentEdge circle[radius=1.5pt];
      \node[rheon/alarm] at \RheonReassignmentAlarm {};
      \draw[rheon/axis] (0,1) -- (0,0) -- (1,0);
      \foreach \y/\value in {0/0,.5/0.5,1/1} {
        \draw[rheon/axis] (0,\y) -- ++(-2.3pt,0);
        \node[rheon/tick,anchor=east,xshift=-5pt] at (0,\y) {\value};
      }
      \foreach \x/\date in \RheonReassignmentTicks {
        \draw[rheon/axis] (\x,0) -- ++(0,-2.3pt);
        \node[rheon/tick,anchor=north,yshift=-5pt] at (\x,0) {\date};
      }
    \end{scope}

    \begin{scope}[shift={(7.39,0)},x=5.89cm,y=2.28cm]
      \fill[KairoPaper] (0,0) rectangle (1,1);
      \fill[KairoAmberTint] (\RheonWaitingOnset,0) rectangle (\RheonWaitingEnd,1);
      \draw[rheon/grid] (0,.5) -- (1,.5) (0,1) -- (1,1);
      \draw[rheon/truth] (\RheonWaitingOnset,0) -- (\RheonWaitingOnset,1)
        (\RheonWaitingEnd,0) -- (\RheonWaitingEnd,1);
      \draw[rheon/threshold] (0,\RheonWaitingThreshold) -- (1,\RheonWaitingThreshold);
      \RheonWaitingPath
      \filldraw[draw=KairoBlue,fill=white,line width=.55pt] \RheonWaitingEdge circle[radius=1.5pt];
      \node[rheon/alarm] at \RheonWaitingAlarm {};
      \draw[rheon/axis] (0,1) -- (0,0) -- (1,0);
      \foreach \x/\date in \RheonWaitingTicks {
        \draw[rheon/axis] (\x,0) -- ++(0,-2.3pt);
        \node[rheon/tick,anchor=north,yshift=-5pt] at (\x,0) {\date};
      }
    \end{scope}

    \draw[rheon/series] (1.02,-.82) -- (1.49,-.82);
    \fill[KairoBlue] (1.255,-.82) circle[radius=1.05pt];
    \node[anchor=west,inner sep=0pt] at (1.61,-.82) {Window score};
    \draw[rheon/threshold] (4.06,-.82) -- (4.53,-.82);
    \node[anchor=west,inner sep=0pt] at (4.65,-.82) {Alarm threshold};
    \fill[KairoAmberTint] (7.36,-.93) rectangle (7.83,-.71);
    \draw[rheon/truth] (7.595,-.93) -- (7.595,-.71);
    \node[anchor=west,inner sep=0pt] at (7.95,-.82) {Annotated drift};
    \node[rheon/alarm] at (10.85,-.82) {};
    \node[anchor=west,inner sep=0pt] at (11.04,-.82) {Estimated alarm};
  \end{tikzpicture}
  \caption{Recent-window Kairo SOM scores for Rheon Single 11 (a) and Rheon Single 15 (b). The annotated change is a point for resource reassignment and an interval for gradual waiting-time change. The diamond marks the alarm at the start of its signal window. Hollow circles mark finite scores from windows excluded from the alarm rule.}
  \label{fig:assessment-rheon-signals}
\end{figure}


\paragraph{Case composition.}
On Rheon Multi 03, the dominant region changed three times and returned to region 4 at the last change. Kairo's inter-case SOM with a previous-window reference matched all three shifts and raised two unmatched alarms. The recent-window version matched two with no unmatched alarms. The core \softwarepm{} PELT indicators and flow-and-timing MMD profiles matched none. We then ran a diagnostic with the case attributes supplied explicitly. PELT used region shares and average case amounts, and MMD used region as a categorical feature and amount as a numeric feature. Each detected all three shifts with zero unmatched alarms, including at the 1\% margin. This diagnostic evaluates an attribute choice made after looking at the scenario. It is reported separately from the core ranking.

\subsection{Sensitivity to the Matching Margin}
\label{sec:assessment-sensitivity}

For \ref{rq:margin}, we took the profiles that led each library at the primary margin in \autoref{tab:assessment-main-results} and scored them at both alternative margins. At 1\%, micro F1 was 0.231 for Kairo and 0.419 for version clustering. At 5\%, their scores were 0.400 and 0.516. Version clustering retained the higher micro F1 at both margins.

With the alarms fixed, the margin changes which alarms can match the annotations. For example, the wide DFG likelihood-ratio profile placed two boundaries in Ostovar Loop about 26 hours after their annotated case starts. Both boundaries matched only at the widest margin.
